\begin{statement}{C.2.RAT.4}{Proved}{Every rational degree-two tacnode cover has positive Euler number}
Retain the torus-fibre, rational-normalization and specified
ruled-boundary morphism of C.2.RAT.1. Suppose $C$ is not
contracted, $H=h(C)_{\rm red}$ has generic transverse type
$A_3$, and the actual map $C\to H^\nu$ has degree two.
Then $e(S)\ge1$. No smoothness or local complete intersection
assumption on the image curve is required.
\end{statement}
\begin{proof}
Write $\Sigma=H^\nu\simeq\PP^1$, $G=\Gamma_{\rm red}$,
$\delta=p_a(H)\le2$, and
$s=\#\nu^{-1}(\operatorname{Sing}H)\le3$ as in
(C.RA.5)--(C.RA.6). There are exactly two generic
normalization branches above $H$, both supplied by the
degree-two map from $C$. Thus no vertical strict conductor
component dominates $H$.

We first justify using the lci jet construction without
assuming that property. The low-delta classification in the
proof of C.2.RAT.1 applies to this same reduced image curve.
Its three-independent-axis germ is excluded by
$b_1(H)\le b_1(G)\le1$. The remaining non-lci possibilities
are the space cusp with semigroup $\langle3,4,5\rangle$
and a cusp with a transverse smooth branch. The
genus-zero double-cover calculation at the end of A.5.18
excludes both: the space-cusp inequalities force
$q_0=-1$ and two simple zeros of the cleared quartic,
one over $H_{\rm reg}$, contradicting (A.T.12); for the
cusp-plus-line ring its cleared quartic has nonpositive
degree and cannot have two branch points. Therefore $H$
is lci at all its points.

Let $\Xi$ be the normalization conductor divisor of $H$,
and use the actual jet lines and sections in
(A.T.5)--(A.T.7). Put
\[
 A=\deg\operatorname{div}(a)=-2\deg Q,\qquad
 \sigma\in H^0(\Sigma,L^2(4\Xi)),\qquad
 \deg\operatorname{div}(\sigma)=8-A.
\tag{C.RA.26}
\]
Here $A$ is a nonnegative even integer and $\sigma$ is
the conductor-cleared version of $\Delta=ae-b^2/4$.
The degree-two cover $C\to\Sigma$ has two branch points,
exactly the odd-order zeros of $\sigma$. Let $B_s$ be
the number of these branch points lying over singular
points of $H$; there are $2-B_s$ over its smooth locus.
Every smooth branch point has zero order at least three
in $\sigma$, by the actual coordinate exclusion
(A.T.10)--(A.T.12).

The section $a$ must vanish at every point of $\Sigma$
over a singular point of $H$. To see this, the Hessian
matrix has entries in $\OO_H$. If its normalized scalar
$a$ is nonzero on one branch above $p$, its primitive
rank-one factor has a unit component there. A diagonal
entry consequently has a nonzero value at $p$. That
same entry has the same nonzero value on every branch,
so $a$ is nonzero at all normalization points above $p$.
In that event (A.T.13) removes the conductor denominator:
$\Delta$ is holomorphic in $L^2$ there. The orders of
$\sigma$ above $p$ are at least those of $4\Xi$,
which are even. Its total zero divisor must additionally
contain the two parity contributions of the two branch
points, whether or not they lie over $p$. Thus
\[
 \deg\operatorname{div}(\sigma)
       \ge4\deg(\Xi|_{\nu^{-1}(p)})+2
       =8\delta_p+2\ge10,
\tag{C.RA.27}
\]
contrary to (C.RA.26). This proves the claimed vanishing.

It follows in particular that
\[
                              A\ge2s-2.
\tag{C.RA.28}
\]
For $s\le1$ this is automatic; for $s=2$ the two
zeros give $A\ge2$. If $s=3$, the bounds
$\delta\le2$ and $b_1(H)\le1$ force exactly two
singular points, each of delta invariant one: one node
and one cusp. The two node preimages contribute at
least one each. At the cusp the ring is
$\mathbb C\{t^2,t^3\}$, so any vanishing Hessian entry
has order at least two. Dividing such a diagonal entry
by the square of a unit component of the primitive
normalized factor shows that $a$ has order at least
two there. Thus $A\ge1+1+2=4=2s-2$.

We now check exactly which smooth image points can
receive two vertical boundary attachments on their two
conductor preimages. At a smooth point with $a\ne0$,
A.5.17 gives the constant $A_3$ product and no vertical
attachment. Suppose instead that $a=0$ and $\Delta$
is a unit. In the actual coordinates (A.T.10), the
coefficient $b$ is then a unit. Substitute $x=y^2\lambda$
in that equation and divide by $y^4$. The result is
holomorphic in $(\lambda,y,z)$ and has, at the point
in question, the simple root
$\lambda=-e(0)/b(0)$. The implicit-function theorem
gives a holomorphic $\lambda(y,z)$ and an actual
factorization
\[
 f_0=F_1F_2,\qquad
 F_1=x-y^2\lambda(y,z),\qquad
 F_2|_{F_1=0}=y^2\,u(y,z),\quad u(0)\ne0.
\tag{C.RA.29}
\]
The last equality is the restriction of
$\partial f_0/\partial x$ to $F_1=0$; its quotient
by $y^2$ has constant value $b(0)$.
The branch $F_1=0$ is smooth, and its conductor in the
whole reduced hypersurface is exactly the divisor
$2\{y=0\}$, by the intersection factor in
(C.RA.22). There is therefore no vertical conductor
near this normalization point. Its minimal resolution
is an isomorphism there as well. Since $\Delta$ is a
unit, this is one of the two distinct conductor
preimages over the given point of $\Sigma$.
Consequently a smooth nonbranch value with two vertical
attachments has a positive even zero order of $\sigma$,
at least two.

Let $d_{\rm reg}$ count distinct smooth values with
two such attachments. The two branch points and these
distinct nonbranch zeros give
\[
             3(2-B_s)+B_s+2d_{\rm reg}\le8-A.
\tag{C.RA.30}
\]
Among the $s$ points above the singular locus, the
$B_s$ branch points each have only one preimage on
$C$. Thus at most $s-B_s$ of these values can have
two attachments. If $d$ counts all values on
$\Sigma$ with two attachments, (C.RA.28)--(C.RA.30)
give
\[
 d\le s-B_s+d_{\rm reg}
       \le s+1-\frac A2\le2.
\tag{C.RA.31}
\]

In fact $d\le1$. If $d=2$, (C.RA.31) requires
$A\le2s-2$. For $s=0$ this is impossible, and for
$s=1$ it contradicts the vanishing of $a$ over the
singular point. For $s=2$, it forces $A=2$ and
the two singular preimages exhaust the zero divisor
of $a$. For $s=3$, it forces $A=4$, exhausted by
the node and cusp orders $1+1+2$ established above.
In both remaining cases $a$ is nowhere zero over
$H_{\rm reg}$. The constant $A_3$ product there
has neither a cover branch point nor a vertical
attachment. Thus both cover branch points lie
over $\operatorname{Sing}H$, and only the other
$s-2\le1$ singular preimages can have two
attachments. This again contradicts $d=2$.
Consequently
\[
                            d\le1.
\tag{C.RA.31a}
\]

It remains to connect this local count to the exact
surface Euler number. As in C.2.RAT.3, at least
$K\ge r-N_0$ original positive vertical boundaries
have a component which is not $\pi$-exceptional.
None maps onto $H$, so they give $K$ distinct points
of $C$ anchored to image curves other than $H$.
After projection to $\Sigma$ their number of distinct
normalization-incidence edges is $k=K-d'$ where
$d'\le d\le1$: a degree-two fibre has at most two
points, and the loss of one occurs precisely at a
value with two selected attachments. Apply the graph
construction of (C.RA.24) to these $k$ distinct
edges. It gives $n_G\ge k$ and $e(G)\ge n_G$.
The actual Euler decomposition therefore yields
\[
\begin{split}
 e(S)&=2-r+N_0+e(G)
       \ge2-K+e(G)\\
     &\ge2-K+k=2-d'\ge1.
\end{split}
\tag{C.RA.32}
\]
Every contraction, covering fibre and graph edge in
this count comes from the specified maps
$Z\to T\to S$ and $C\to\Sigma\to H$.
\end{proof}


% C.2.UN.67 from part_C2_rational_hopf_interface.tex

